%========================================================================================
\section{Gestion des voyages, destinations, activités et \mbox{participants}}
\sectionmark{Voyages et participants}
\label{sec:gestion}
%========================================================================================

\subsection{QuerySets réutilisables}

Les requêtes fréquentes sont regroupées dans un \code{QuerySet} personnalisé, exposé par \code{Trip.objects} (\code{TripQuerySet.as_manager()}). Elles se combinent par chaînage.

\begin{table}[H]
  \centering
  \small
  \caption{Méthodes de \texttt{TripQuerySet}.}
  \label{tab:querysets}
  \begin{tabularx}{\linewidth}{@{}l X@{}}
    \toprule
    \textbf{Méthode} & \textbf{Résultat} \\
    \midrule
    \code{visible_to(user)} & Voyages dont l'utilisateur est propriétaire ou participant, quel que soit son rôle ; vide pour un anonyme \\
    \code{editable_by(user)} & Voyages que l'utilisateur peut modifier : propriétaire, rôle \code{owner} ou \code{editor} \\
    \code{upcoming()} & Voyages à venir, hors voyages annulés ou terminés, triés par date \\
    \code{with_totals()} & Annotations \code{destination_count}, \code{activity_count} et \code{total_cost} \\
    \bottomrule
  \end{tabularx}
\end{table}

La méthode \code{visible_to()} est le point de départ de toute vue et de tout resolver GraphQL. Elle utilise une sous-requête \code{Exists()} plutôt qu'une jointure sur les participants. Une jointure dupliquerait chaque voyage autant de fois qu'il a de participants, et fausserait les sommes de \code{with_totals()}.

\begin{lstlisting}[language=Python,caption={Voyages visibles et totaux (\texttt{apps/trips/models.py}, extrait).},label={lst:visible}]
def visible_to(self, user):
    if user is None or not user.is_authenticated:
        return self.none()
    # Exists() plutôt qu'une jointure : aucun doublon, et with_totals() reste exact.
    membership = TripParticipant.objects.filter(trip=OuterRef("pk"), user=user)
    return self.filter(Q(owner=user) | Exists(membership))

def with_totals(self):
    return self.annotate(
        destination_count=Count("destinations", distinct=True),
        activity_count=Count("destinations__activities", distinct=True),
        total_cost=Coalesce(
            Sum("destinations__activities__cost"),
            Value(0),
            output_field=DecimalField(max_digits=12, decimal_places=2),
        ),
    )
\end{lstlisting}

Le calcul de \code{total_cost} suit la chaîne voyage $\rightarrow$ destinations $\rightarrow$ activités. Chaque activité n'apparaît qu'une fois dans cette jointure : la somme est exacte. Le code interdit d'y ajouter un \code{Count()} sur une seconde relation, car cette jointure supplémentaire multiplierait les lignes. Pour la même raison, le tri de la liste des voyages ne propose aucun champ de relation.

\subsection{Relations inverses et préchargement}

Les relations inverses du tableau~\ref{tab:relations} évitent d'écrire des requêtes explicites : \code{trip.memberships.select_related("user")} liste les participants, \code{destination.activities.all()} les activités d'une étape. Le filtrage à travers les relations reste possible dans les deux sens, par exemple \code{Activity.objects.filter(destination__trip=trip)}.

Le préchargement empêche le problème des \og{}N+1 requêtes\fg{}. La liste des voyages charge le propriétaire par \code{select_related("owner")}. Le détail d'un voyage précharge ses destinations, triées par position, puis leurs activités, à l'aide d'un objet \code{Prefetch}. Le nombre de requêtes SQL ne dépend donc pas du nombre de destinations.

\subsection{Agrégations}

L'action \code{GET /api/trips/{id}/stats/} calcule le budget restant et la répartition des coûts par catégorie. Tous les calculs sont confiés à PostgreSQL : \code{aggregate()} pour les totaux, \code{values().annotate()} pour le regroupement.

\begin{lstlisting}[language=Python,caption={Statistiques par catégorie (\texttt{apps/trips/services.py}, extrait).},label={lst:stats}]
activities = Activity.objects.filter(destination__trip=trip)
totals = activities.aggregate(count=Count("pk"), total=Sum("cost"), average=Avg("cost"))
by_category = (
    activities.values("category")
    .annotate(count=Count("pk"), total_cost=Sum("cost"), average_cost=Avg("cost"))
    # order_by() explicite : l'ordre par défaut d'Activity ajouterait date au GROUP BY.
    .order_by("-total_cost", "category")
)
\end{lstlisting}

Le commentaire de l'extrait~\ref{lst:stats} signale un piège de l'ORM. Le tri par défaut d'\code{Activity} porte sur la date. S'il était conservé, Django ajouterait la date au \code{GROUP BY}, et chaque catégorie serait éclatée par jour.

\subsection{Services métier}

\paragraph{Création.} Le service \code{create_trip()} valide le voyage (\code{full_clean()}), l'enregistre, puis inscrit le propriétaire comme participant \code{owner}, dans une seule transaction. Un signal \code{post_save} aurait produit le même effet, mais de façon implicite. Le service rend l'effet de bord visible à l'appel.

\paragraph{Clonage.} Le service \code{clone_trip()} copie un voyage, ses destinations et leurs activités dans un nouveau brouillon. Il utilise \code{bulk_create()}, qui insère chaque niveau en une requête. Sous PostgreSQL, \code{bulk_create()} renvoie les clés primaires créées : les destinations copiées servent directement de parents aux activités copiées.

\begin{lstlisting}[language=Python,caption={Clonage transactionnel (\texttt{apps/trips/services.py}, extrait).},label={lst:clone}]
@transaction.atomic
def clone_trip(source: Trip, *, owner, title=None) -> Trip:
    trip = create_trip(owner=owner, title=title or f"Copie de {source.title}", ...,
                       status=Trip.Status.DRAFT)
    destinations = list(source.destinations.prefetch_related("activities"))
    # bulk_create renvoie les clés primaires sous PostgreSQL : les copies servent de parents.
    copies = Destination.objects.bulk_create(_copy(d, trip=trip) for d in destinations)
    Activity.objects.bulk_create(
        _copy(activity, destination=copy)
        for original, copy in zip(destinations, copies, strict=True)
        for activity in original.activities.all()
    )
    return trip
\end{lstlisting}

\paragraph{Participants.} Le propriétaire ajoute un participant par son courriel, avec le rôle \code{editor} ou \code{viewer}. Le rôle \code{owner} ne s'attribue pas à l'invitation : il se transfère ensuite par un changement de rôle. L'ancien propriétaire devient alors éditeur. Le propriétaire ne peut être ni retiré ni rétrogradé directement ; le voyage resterait sans propriétaire. Un ajout en double viole la contrainte \code{trips_participant_unique_trip_user}. Le service laisse la base trancher, ce qui élimine toute course entre deux ajouts simultanés.

\subsection{Administration et données de démonstration}

L'administration de Django (\code{/admin/}) affiche les voyages avec leurs participants et leurs destinations en ligne (figure~\ref{fig:admin}). La commande \code{seed_demo} crée trois comptes et quatre voyages (Montréal, Paris, Tokyo, Dakar) qui couvrent chaque rôle. Elle est idempotente, et ses dates sont relatives au jour de l'exécution : un voyage terminé, un en cours, un planifié et un brouillon.

\capture{admin.png}{Administration Django : fiche d'un voyage, avec ses participants et ses destinations en ligne.}{fig:admin}
